\chapter{Blocks 5--6: R, Bioconductor \& Differential Expression}

\textbf{Primary:} Bioconductor, differential expression, ggplot2\quad|\quad\textbf{Interleave:} Statistics concepts, Unix file prep

\begin{brainbox}[Why R when you just learned Python?]
R has Bioconductor --- thousands of biology-specific packages with nothing comparable in Python. DESeq2, edgeR, GenomicRanges, Seurat all live here. You need both languages.
\end{brainbox}


\section{Block 5: R Basics + Bioconductor}

\sprint{1}{R in 25 Minutes --- Vectors, Lists, Data Frames}{25}
\begin{itemize}[nosep]
  \item \textbf{Vectors.} \texttt{x <- c(1, 2, 3); mean(x)}. Everything in R is a vector --- even a single number is a length-1 vector. Vectors are \textit{homogeneous}: one type throughout (numeric, character, logical).
  \item \textbf{Lists.} Heterogeneous containers:
    \begin{lstlisting}
res <- list(pval = 0.003,
            estimate = 1.42,
            samples = c("A", "B", "C"))
res$pval          # 0.003
res[["samples"]]  # c("A", "B", "C")
    \end{lstlisting}
    Access by name with \texttt{\$} or \texttt{[[}. Every complex R object --- \texttt{lm()} output, a \texttt{DESeqDataSet}, a phylo tree --- is a list underneath. Learn \texttt{str(obj)} now; you'll use it constantly to dig into unfamiliar objects.
  \item \textbf{Data frames.} A list of equal-length columns with row-and-column indexing: \texttt{df <- read.csv('samples.tsv', sep='\textbackslash t')}. First three commands on any new data frame: \texttt{head(df)}, \texttt{names(df)}, \texttt{str(df)}.
  \item \textbf{Indexing gotchas.} R is 1-based: \texttt{df[1,]} is row 1, not row 2. Negative indices \textit{exclude}: \texttt{df[-1,]} drops row 1. Logical vectors filter: \texttt{df[df\$pval < 0.05, ]}.
\end{itemize}
\begin{winbox}[Checkpoint]
You loaded a TSV, inspected its structure with \texttt{str()}, pulled one column and one row, and know the difference between a vector, a list, and a data frame. Apply functions and the tidyverse come next.
\end{winbox}

\sprint{2}{Lists, the Apply Family \& Functional R}{25}
R grew up as a functional language: instead of writing for-loops, you hand a function to another function that runs it over a collection. This pattern is everywhere in R code --- you need to read it fluently even if you'd rather write loops.
\begin{itemize}[nosep]
  \item \textbf{\texttt{lapply()} --- run f over a list, return a list.}
    \begin{lstlisting}
genes <- list(BRCA1 = counts1, BRCA2 = counts2)
means <- lapply(genes, mean)
# means$BRCA1, means$BRCA2
    \end{lstlisting}
  \item \textbf{\texttt{sapply()}} does the same and simplifies to a vector/matrix when it can. Convenient interactively; prefer \texttt{vapply()} in scripts --- it type-checks the return shape and fails loudly when a single element deviates.
  \item \textbf{\texttt{apply()} --- walk rows or columns of a matrix.}
    \begin{lstlisting}
row_means <- apply(counts_matrix, 1, mean)
col_sds   <- apply(counts_matrix, 2, sd)
    \end{lstlisting}
    First arg = matrix; \texttt{1} = rows, \texttt{2} = columns; third arg = the function.
  \item \textbf{\texttt{mapply()} / \texttt{Map()} --- over multiple parallel inputs.} \texttt{mapply(t.test, group\_a\_list, group\_b\_list)} runs a t-test pairing element 1 of each, element 2 of each, etc.
  \item \textbf{\texttt{tapply()} --- split by a grouping vector, then apply.} \texttt{tapply(expression, tissue, mean)} is ``mean expression per tissue'' in one call.
  \item \textbf{Why this matters.} DESeq2 results come back as a list; \texttt{apply()} is how you compute per-gene statistics; \texttt{lapply()} appears in every production R pipeline. Loops are not wrong, but apply-family code is shorter and more idiomatic.
  \item \textbf{The modern equivalent.} The \texttt{purrr} package (part of the tidyverse) gives you \texttt{map()}, \texttt{map\_dbl()}, \texttt{map\_chr()}, \texttt{map2()}, \texttt{pmap()} --- same idea, stricter return types. We'll use \texttt{purrr::map\_dbl()} in the next sprint.
\end{itemize}
\begin{winbox}[Checkpoint]
You computed per-gene means and per-sample standard deviations on a count matrix without writing a single explicit loop. This is the R way.
\end{winbox}

\sprint{3}{Tidyverse as a Modern Layer}{25}
Base R works, but routine data wrangling is verbose. The tidyverse is a family of packages (\texttt{dplyr}, \texttt{tidyr}, \texttt{readr}, \texttt{purrr}, \texttt{ggplot2}, \texttt{stringr}, \texttt{forcats}, \ldots) that share a consistent grammar: pipe a data frame through a sequence of \textit{verbs}, each returning a modified data frame. Once you internalize the verbs, most analyses become one-glance readable.
\begin{itemize}[nosep]
  \item \textbf{The pipe.} \texttt{\%>\%} (from \texttt{magrittr}, re-exported by dplyr) and the newer base-R \texttt{|>} (R $\geq$ 4.1) both do the same thing: pass the left-hand side as the first argument of the right-hand side. They're Unix \texttt{|} for data frames.
  \item \textbf{The five verbs you'll use constantly.}
    \begin{lstlisting}
library(dplyr)
de_results |>
  filter(padj < 0.05) |>
  mutate(direction = if_else(log2FC > 0,
                             "up", "down")) |>
  group_by(chromosome, direction) |>
  summarise(n = n(), .groups = "drop") |>
  arrange(desc(n))
    \end{lstlisting}
    \texttt{filter} picks rows; \texttt{mutate} adds/changes columns; \texttt{group\_by} + \texttt{summarise} compute per-group statistics; \texttt{arrange} sorts; \texttt{select} picks columns. Five verbs cover 80\% of real analysis.
  \item \textbf{Tibbles.} \texttt{tibble()} creates a modern data frame that prints neatly (no auto-factor conversion, no surprising type coercion). \texttt{readr::read\_csv()} returns a tibble; it behaves like a \texttt{data.frame} everywhere else.
  \item \textbf{Reshaping.} \texttt{pivot\_longer()} and \texttt{pivot\_wider()} (\texttt{tidyr}) go between wide and long formats --- critical for plotting in ggplot2 and for joining disparate datasets.
  \item \textbf{Joins.} \texttt{inner\_join()}, \texttt{left\_join()}, \texttt{anti\_join()}. Merge GWAS summary stats with gene annotations by SNP ID in one line.
  \item \textbf{\texttt{across()} for column-wise operations.} \texttt{mutate(across(starts\_with("sample\_"), log1p))} log-transforms every sample column at once.
\end{itemize}
\begin{tipbox}[When to stay in base R]
Tidyverse is great for interactive analysis and scripts humans will read. Package authors often prefer base R internally --- it has zero dependencies and a more stable API. Both styles are valid; most real projects mix them.
\end{tipbox}
\begin{winbox}[Checkpoint]
You chained five dplyr verbs to go from raw DE results to a ranked table of up/down counts per chromosome. Compare the equivalent base-R code (\texttt{aggregate}, nested \texttt{[}) --- same output, very different ergonomics.
\end{winbox}

\sprint{4}{Bioconductor: GRanges = Genomic Intervals}{25}
\begin{itemize}[nosep]
  \item Install: \texttt{BiocManager::install('GenomicRanges')}
  \item GRanges: \texttt{gr <- GRanges(seqnames='chr1', ranges=IRanges(start=100, end=200))}
  \item \texttt{findOverlaps(peaks, genes)} --- which peaks overlap which genes? This replaces \texttt{bedtools intersect}.
  \item \texttt{rtracklayer::import('file.bed')} --- load BED, GFF, BigWig directly into R
\end{itemize}
\begin{winbox}[Checkpoint]
You performed a genomic interval overlap in R. Same concept as bedtools, different syntax.
\end{winbox}

\sprint{5}{SummarizedExperiment: The Container That Holds Everything}{25}
\begin{itemize}[nosep]
  \item SummarizedExperiment = counts matrix + gene info (rowData) + sample info (colData)
  \item Every Bioconductor analysis tool expects data in this format. Learn it once, use it everywhere.
  \item Load a real dataset: recount3 or GEOquery to pull RNA-seq data directly from public repositories
  \item Explore: \texttt{assay(se)[1:5, 1:5]} --- the first 5 genes and 5 samples
\end{itemize}
\begin{winbox}[Checkpoint]
You loaded a real RNA-seq dataset from a public database into R. Data is ready for analysis.
\end{winbox}


\section{Block 6: Differential Expression + ggplot2}

\sprint{6}{DESeq2: The Gold Standard DE Analysis}{50}
\begin{itemize}[nosep]
  \item Build a DESeqDataSet from your SummarizedExperiment + a design formula (\texttt{\textasciitilde condition})
  \item Run: \texttt{dds <- DESeq(dds); res <- results(dds)} --- the entire analysis in two lines
  \item Understand what happened: normalization (size factors), dispersion estimation (empirical Bayes shrinkage), NB GLM fitting, Wald test
  \item Extract results: \texttt{res[order(res\$padj), ]} --- genes sorted by significance
  \item Make a volcano plot: plot $-\log_{10}(\text{pvalue})$ vs.\ $\log_2(\text{FoldChange})$ with ggplot2
\end{itemize}
\begin{winbox}[Checkpoint]
You ran a full differential expression analysis on real data. This is publishable work.
\end{winbox}

\sprint{7}{ggplot2: Grammar of Graphics}{25}
\begin{itemize}[nosep]
  \item The formula: \texttt{ggplot(data, aes(x, y)) + geom\_point() + theme\_minimal()}
  \item Layers: \texttt{+ geom\_violin() + geom\_boxplot(width=0.1)} --- combine geoms for rich plots
  \item Facets: \texttt{+ facet\_wrap(\textasciitilde chromosome)} --- split by a variable automatically
  \item Save: \texttt{ggsave('figure1.pdf', width=8, height=6)} --- publication-ready output
\end{itemize}
\begin{winbox}[Checkpoint]
You created a multi-panel publication figure. This goes in your portfolio.
\end{winbox}

\sprint{8}{Compare DE Methods (Interleaving!)}{25}
\begin{itemize}[nosep]
  \item Run the same dataset through limma-voom: \texttt{voom(counts) + lmFit + eBayes + topTable}
  \item Compare results: how many genes are significant in DESeq2 but not limma? Vice versa?
  \item Make a Venn diagram or UpSet plot of the overlap --- this is how real papers compare methods
  \item Think about why: DESeq2 and limma make different statistical assumptions. Which matter for your data?
\end{itemize}
\begin{winbox}[Checkpoint]
You compared two statistical methods on the same data and thought critically about why results differ.
\end{winbox}

\begin{restbox}[End of Blocks 5--6]
You now have working fluency in bash, Python, and R, and can run a differential expression analysis. Take a break before moving on.
\end{restbox}

\subsection{Checkpoint:}
\begin{itemize}[nosep]
  \item You can manipulate data in R with tidyverse
  \item You understand GRanges and SummarizedExperiment
  \item You can run a DESeq2 analysis end-to-end
  \item You can make publication-quality ggplot2 figures
  \item You can explain what empirical Bayes shrinkage does and why DESeq2 uses it
\end{itemize}

\subsection{Resources}
\begin{itemize}[nosep]
  \item R for Data Science, 2nd Ed (r4ds.hadley.nz) --- free online, excellent
  \item DESeq2 vignette (follow the Beginner's Guide first, then the full vignette)
  \item StatQuest YouTube: the DESeq2 and edgeR episodes are outstanding
  \item Bioconductor workflows: bioconductor.org/packages/release/workflows/
\end{itemize}
